\section{Joint-Design Parameter Studies}\label{sec:joint-study}
The new computations ask how the optimal continuous book changes with the promise, realization ceiling, and selected-level charges, and how a finite grid approaches that optimum. They are synthetic studies of the declared model, not an empirical calibration. Every optimum in the first two studies is obtained by the global rational face search of Theorem~\ref{thm:continuous}; its reconstructed book is then checked by the independent fixed-book multiplier certificate. This is stronger than treating a locally optimized book as a continuous benchmark, but it is not a substitute for the theorem's proof.

\paragraph{The institutional gap across promises.}
Take caps $(1/4,1/2,3/4)$ with equal probabilities, $r=2$, $q=1$, $\gamma_j=1$, $m=2$, and no codeword charges. Table~\ref{tab:r42-gap} reports eight promises and compares the unrestricted expected, pathwise, and deterministic optima. The expected/pathwise gap is zero at each of the five sampled promises through $5/12$, and strictly positive at each sampled promise from $11/24$ to saturation. At $B=11/24$, the exact gap is $11/192$ even though the old transport certificate gives zero. Thus the sufficient interval $11/24<B\leq1/2$ is conservative; it is not the maximal separation interval. These samples alone do not identify the exact transition point. Pathwise and deterministic values coincide along this particular uncharged sample path, not universally in the joint charged model: Proposition~\ref{prop:pathwise-interior} gives an exact counterexample to that extrapolation.

\paragraph{Risk and actual level prices.}
For caps $(1/4,1/2,3/4)$ and probabilities $(7/20,3/5,1/20)$, retain the same quadratic primitives and $m=2$. We cross four promises $B\in\{1/3,3/8,2/5,17/40\}$, four ceilings $\delta\in\{0,1/16,1/8,1\}$, and two price schedules: zero and $\rho(c)=c/50+1/200$. This gives 32 joint continuous designs. Table~\ref{tab:r42-joint} displays the $B=2/5$ and saturated slices; the archived rational records include every target, intermediate tier, lottery, and enabled level. Because charges enter before optimization, changes in the selected book are not post hoc accounting adjustments. The algorithm permits off-cap levels, endogenous branch targets, and level collisions throughout these comparisons.

\paragraph{Continuous benchmark versus certified grids.}
At $B=2/5$, $\delta=1/16$, and $\rho(c)=c/50+1/200$, the continuous net optimum is $7169/10400$. Table~\ref{tab:r42-mesh} compares it with uniform grids of 9 through 129 points. The realized gap falls from $19/10400$ to $57/166400$ along these grids. The theorem's worst-case bound is deliberately conservative: here $K_2=493/200$, whereas the observed errors are much smaller than $K_2h$. The unchanged optimum over some successive grids is not a violation of convergence; nested grids give monotone values but need not improve at every refinement. All displayed grid policies satisfy the original promise and realization ceiling exactly.

\paragraph{Scaling of the new problem.}
Table~\ref{tab:r42-scaling} reports separate measurements for all-promise charged catalog enumeration, rather than reusing the earlier 16,384-branch saturated benchmark. We use $m=2$, caps $b_j=1/4+j/[2(k+2)]$, equal probabilities, $\gamma_j=1+((j-1)\bmod4)$, $B=3\bar b/4$, $\delta=1/16$, and the preceding affine charge. These small exact-arithmetic runs expose the dependence on both branch count and catalog size; they do not establish large-scale performance of the exponential continuous face algorithm. The continuous reference runs have three branches, except for the one-branch pathwise example and the two-branch regression cases. Each timing is one measured solve, excluding import and reader compilation, with no averaging or hardware-independent interpretation. The environment and machine-readable timings are retained separately from exact values.
